% Page 108

\seite{108}

\margmark{5. April}%
Das neue Schuljahr hat heute begonnen. Neu auf\ohb
genommen wurden 5 Kinder: 2 Knaben und 3 Mäd\ohb
chen. Die Schule besitzt jetzt 50 Kinder. Wegen\ob
Raummangel müßte wohl\unsicher weiterhin eine neue Bank\ob
3 sitzig beschafft werden.\ob
Mit Wirkung vom 1.\ April 1932 ist Herr Lehrerwart\ob
Kesseler endgültig nach Bitburg versetzt worden. Die\ob
Lehrerstelle Kreis Neuendorf wird mit dem 14.\unsicher\ob
aufgeben.

\margmark{10. April}%
Bei der heute stattgefundenen 2.\ Wahlgang\ob
zur Reichspräsidentenwahl wurden hier 113 Stimmen\ob
abgegeben. Davon fielen auf von Hindenburg 88, auf\ob
Hitler 20 Stimmen. Ungültige waren 5 Stimmen.\ob
von Hindenburg ist wiederum mit großer Mehrheit\ob
zum Reichspräsidenten gewählt worden.

\margmark{25. April}%
Gestern fand die Wahl zum preußischen Landtag\ob
statt. Von den 149 Wahlberechtigten haben 128 ge\ohb
wählt.\unsicher{} Wahlrecht Gebung erreicht. Es erhielten Eifeler\ob
Volkspartei nationale\unsicher{} 30, Liste 2a. Nationale\ob
Sammlung Gust. Anders 2, Liste 3: Preußische Zentrumspartei\unsicher\ob
\unsicher 73, Liste 6: Deutscher Landvolk 3, und Liste 8: National\ohb
sozialistsche Deutsche Arbeiterpartei\unsicher{} 59\unsicher{} 13 Stimmen.\ob
2 Stimmen waren ungültig.

Nach der Bearbeitung der schweren Schäden in Seffern\ob
fand seit Beginn des Schuljahres bis auf weiteres\unsicher\ob
ein eine sonntäglich\unsicher{} Volksabstimmung statt.

\margmark{14. Mai}%
Die dreitägigen Pfingstferien dauern vom\ob
14.\ bis 23.\ Mai einschließlich.

\margmark{21. Juni}%
Mit dem heutigen Tag beginnen die Herbstferien.\ob
Die 14 Tage dauern.\ob
Die Fleischbeschauungstelle hat des Letzten Ehren\ohb
berichte\unsicher{} worden fortwährend\unsicher{} vernachlässigt\ob
der Landwirt hat Fleisch\unsicher{} mal einmal Fleisch durch\unsicher\ob
\unsicher vorschrif\unsicher{} die Bauern. Die Wolfern\unsicher
\pend
